\section{Where the active ingredients may lie}\label{sec:ingredients}
\subsection{Component and dismantling evidence}
The foundational dismantling trial (depression, $N=150$) found no evidence that the complete cognitive treatment outperformed behavioural activation alone or behavioural activation plus automatic-thought skills at termination or at 6 months \citep{jacobson1996component}, and a large non-inferiority trial found behavioural activation as effective as CBT (mean PHQ-9 difference $0.1$; \citealt{richards2016cobra}). In anxiety:
\begin{itemize}
\item In a 108-trial meta-analysis, cognitive therapy, exposure and their combinations were efficacious with ``no differential efficacy for any treatment components'' \citep{norton2007metaanalytic}.
\item In a full-factorial trial of internet CBT for social anxiety ($N=464$), conditions with psychoeducation and exposure were more effective, while conditions with cognitive restructuring or attention training were not superior at post-treatment; the advantage of exposure faded at follow-up \citep{sipka2025factorial}.
\item Adding cognitive restructuring to prolonged exposure did not enhance outcome in assault survivors with PTSD \citep{foa2005pecr}, nor in children with social anxiety \citep{demooij2023social}.
\item A component network meta-analysis in panic disorder (72 studies) singled out interoceptive exposure, not cognitive restructuring, as efficacy-enhancing \citep{pompoli2018dismantling}.
\item In GAD, the network meta-analysis authors concluded that pure cognitive restructuring was not a decisively active ingredient when considered alone \citep{papola2024gad}; the paper was later corrected for a duplicated cohort and an error in standard-error handling \citep{papola2024correction}.
\end{itemize}

\subsection{Evidence against the claim}
These findings are not uniform.
\begin{itemize}
\item Cognitive therapy beat exposure in social phobia in a meta-analysis of 20 direct comparisons ($N=1{,}308$), with no difference in PTSD, OCD or panic \citep{ougrin2011exposure}.
\item In a trauma trial, adding cognitive restructuring to exposure was associated with PTSD rates of 31\% versus 63--75\% without it at 6 months \citep{bryant2008cr}.
\item In-session cognitive restructuring correlates with outcome ($r=.35$), though from only four studies and correlationally \citep{ezawa2023restructuring}.
\item In a single-session trial ($N=249$), restructuring before exposure worked about as well as behavioural experiments and better than exposure without explicit processing \citep{johnco2025restructuring}.
\item The exposure--cognition split is blurred: exposure may itself work by changing the expectancy of harm \citep{hofmann2008extinction}.
\end{itemize}

\subsection{Why the component evidence is weak on both sides}
A review of 16 component studies concluded that they lack ``the statistical power nor the quality to draw any meaningful conclusion,'' because none except one could detect an effect smaller than $g=0.55$ \citep{cuijpers2017components}. Mediation evidence is correlational: cognitive change tracks symptom change however it is produced \citep{cristea2015dysfunctional,lorenzoluaces2015complicated}, cognitive mediation of symptom reduction was limited in a depression cohort \citep[median $r=.06$;][]{vittengl2014mediators}, and reviewers say no common or specific factor meets causal criteria \citep{kazdin2007mediators,cuijpers2019commonfactors}. A null component result is thus weak evidence that restructuring adds nothing, and a positive correlation is weak evidence that it matters.

\subsection{The strongest defence: practice when calm}
A defender can answer that patients do not restructure during a panic; they practise when calm so the response becomes automatic. We accept the force of this. It changes the claim, though: if rehearsal during calm periods does the work, the active ingredient is practice and habit formation, and the relevant question is the executive cost of each rehearsal. Exposure, which is central to anxiety CBT, is behavioural, and optimised exposure protocols are explicitly framed around expectancy violation and extinction \citep{craske2014maximizing}.
